\documentclass{article}
\usepackage{amsmath}
\usepackage{amsthm}

\newtheorem{theorem}{Problem}

\begin{document}

\begin{theorem}
What is $\sqrt{53}$ in simplest radical form?
\end{theorem}

\begin{proof}
Step 0: Step 1. I need to simplify the square root of 53, which is written as √53.
Step 1: Step 2. To simplify a square root, I look for perfect square factors of the number under the radical.
Step 2: Step 3. Perfect squares are numbers like 1, 4, 9, 16, 25, 36, 49, and so on, because they are squares of integers.
Step 3: Step 4. I need to check if 53 has any perfect square factors other than 1.
Step 4: Step 5. First, I check if 53 is a prime number. A prime number has no positive divisors other than 1 and itself.
Step 5: Step 6. I test divisibility of 53 by prime numbers less than its square root. The square root of 53 is approximately 7.28, so I test primes up to 7.
Step 6: Step 7. Testing divisibility by 2: 53 is odd, so it's not divisible by 2.
Step 7: Step 8. Testing divisibility by 3: Sum of digits of 53 is 5 + 3 = 8, which is not divisible by 3, so 53 is not divisible by 3.
Step 8: Step 9. Testing divisibility by 5: 53 does not end with 0 or 5, so it's not divisible by 5.
Step 9: Step 10. Testing divisibility by 7: 7 times 7 is 49, and 7 times 8 is 56, which is greater than 53. So, 53 is not divisible by 7.
Step 10: Step 11. Since 53 is not divisible by any of these primes, it has no divisors other than 1 and itself. Therefore, 53 is a prime number.
Step 11: Step 12. Because 53 is prime, its only factors are 1 and 53. There are no perfect square factors other than 1.
Step 12: Step 13. Therefore, √53 cannot be simplified further.
Step 13: Step 14. The simplest radical form of √53 is √53 itself.

        # Answer

        \sqrt{53}
\end{proof}

\end{document}
